% TOP_PROBLEM: {"id": 3197, "title": "Packing T-joins", "category_id": 3, "category": {"id": 3, "name": "graph_theory", "display_name": "Graph Theory", "description": "Problems involving graphs, networks, and their properties.", "slug": "graph-theory", "order_index": 3, "created_at": "2026-07-31T15:26:25.671Z"}, "status": "open", "problem_number": "OPG-144", "set_id": 12}
% FIRST_POSTED: 2026-10-01
% ENTRY_KIND: statement_only
% UnsolvedMath ID 3197; source catalogue code OPG-144
% !TeX program = lualatex
% Definition and sources only; this file is not an LLM solution attempt.
\documentclass[11pt,a4paper]{article}
\usepackage[margin=25mm]{geometry}
\usepackage{fontspec}
\setmainfont{lmroman10-regular.otf}[BoldFont=lmroman10-bold.otf,
 ItalicFont=lmroman10-italic.otf,BoldItalicFont=lmroman10-bolditalic.otf]
\usepackage{amsmath,amssymb,mathtools}
\usepackage{unicode-math}
\setmathfont{latinmodern-math.otf}
\AtBeginDocument{\let\setminus\smallsetminus}
\usepackage{xurl,hyperref}
\hypersetup{colorlinks=true,urlcolor=blue}
\setcounter{secnumdepth}{0}
\setlength{\parindent}{0pt}
\setlength{\parskip}{5pt}
\setlength{\emergencystretch}{3em}
\begin{document}
\section*{Packing T-joins}\label{3197}
{\small UnsolvedMath ID 3197 · OPG-144 · Graph Theory · OpenGarden}
\subsection{Definitions and mathematical statement}
Let $G=(V,E)$ be a finite loopless undirected multigraph and let $T\subseteq V$ be a nonempty set of even cardinality. A $T$-join is a subset $J\subseteq E$ such that the vertices having odd degree in the spanning subgraph $(V,J)$ are exactly the vertices in $T$. Different edges of a parallel class are distinct available edges. A packing of $T$-joins is a collection of pairwise edge-disjoint $T$-joins; its size is the number of joins, not their total edge count.

For $X\subseteq V$, the edge cut $\delta(X)$ consists of edges with one endpoint in $X$ and the other outside. It is a $T$-cut when $|X\cap T|$ is odd. The conjecture asks whether an absolute constant $C\ge0$ exists such that, for every integer $k\ge1$ and every pair $(G,T)$ for which all $T$-cuts have at least $k$ edges, there is a packing of at least
\[
                         \frac{2k}{3}-C
\]
$T$-joins. Since the packing size is an integer, the displayed lower bound is interpreted literally, or equivalently with its ceiling.

The constant is independent of the graph, the terminal set, and $k$. The positive cut lower bound ensures that each connected component contains an even number of terminals, so the existence of a single $T$-join is not obstructed by disconnectedness. The target is a uniform additive error, not an error allowed to grow with $k$.
\subsection{Short English statement}
Does a minimum T-cut of size k guarantee at least two-thirds of k edge-disjoint T-joins, up to one universal additive constant?
\subsection{Sources}
\begin{itemize}
\item[S1] ULAM.AI and the UnsolvedMath Contributors, supplied problems.json, record 3197 (OPG-144), OpenGarden. Catalogue identity and source transcription; CC BY 4.0.
\url{https://huggingface.co/datasets/ulamai/UnsolvedMath}.
\item[S2] Open Problem Garden, Packing T-joins. Original problem and discussion.
\url{http://www.openproblemgarden.org/op/packing_t_joins}.
\end{itemize}
\end{document}
